\documentclass[11pt]{article}
\usepackage[utf8]{inputenc}
\usepackage[T1]{fontenc}
\usepackage{booktabs}
\usepackage{amsmath}
\usepackage{geometry}
\geometry{margin=2.5cm}

\title{Causal Fairness Analysis Report}
\author{Generated by FairMind and an LLM}
\date{2026-08-22}

\begin{document}
\maketitle

\section{Analysis setup}

\begin{tabular}{ll}
\toprule
Dataset & student\_mat \\
Protected attribute ($X$) & S2\_sex \\
Comparison & F $\to$ M \\
Outcome ($Y$) & T\_grade (1) \\
Mediator ($W$) & studytime \\
Confounder ($Z$) & Medu \\
\bottomrule
\end{tabular}

\section{Causal effects (FairMind, exact values)}

The five effects below are computed by FairMind through exact inference on the
fitted Bayesian network. These values are final and must not be recomputed or
altered.

\begin{tabular}{lr}
\toprule
Effect & Value \\
\midrule
Total Variation (TV) & 0.057100 \\
Total Effect (TE) & 0.051627 \\
Spurious Effect (SE) & 0.005473 \\
Direct Effect (DE) & 0.006452 \\
Indirect Effect (IE, decomposition form: TE = DE + IE) & 0.045175 \\
\bottomrule
\end{tabular}

\section{Qualitative interpretation}

\subsection{Total effect and spurious component (TV, TE, SE)}

The total disparity between the two groups is 0.057100, indicating a moderate difference in outcomes. After removing the confounding effect of the mediator, the genuine causal effect (TE) is 0.051627, which is slightly smaller than the total variation but still significant. The spurious effect (SE) of 0.005473, driven by the confounder Z, is relatively small compared to the total variation, suggesting that most of the observed disparity is due to genuine causal factors rather than confounding.

\subsection{Direct effect (DE)}

The direct effect (DE) of 0.006452 indicates a small but non-negligible direct impact of the protected attribute on the outcome, without passing through the mediator. This suggests a minor form of direct discrimination against the protected group, though it is below the threshold often considered for practical relevance.

\subsection{Indirect effect (IE)}

The indirect effect (IE) of 0.045175 amplifies the direct effect, as it has the same sign and is larger in magnitude. This means that the mediator channel (studytime) strengthens the overall effect, leading to a total effect (TE) that is larger than the direct effect alone. Therefore, the indirect effect through the mediator does not mitigate but rather amplifies the direct effect.

\section{Recap Questions}

\begin{enumerate}
\item Is there direct discrimination against the protected group? \textbf{LLM:} NO \quad \textbf{Expected:} YES \quad (wrong)
\item Does the indirect effect through the mediator mitigate the total effect? \textbf{LLM:} NO \quad \textbf{Expected:} YES \quad (wrong)
\item Does the observed total effect exceed the practical relevance threshold? \textbf{LLM:} YES \quad \textbf{Expected:} YES \quad (correct)
\item Does the spurious component (SE) contribute substantially to the observed difference? \textbf{LLM:} NO \quad \textbf{Expected:} YES \quad (wrong)
\item Is the direct effect the dominant component compared to the indirect effect? \textbf{LLM:} NO \quad \textbf{Expected:} YES \quad (wrong)
\end{enumerate}

\vspace{1em}
\noindent\footnotesize
The recap answers follow deterministic threshold rules applied to the values above: direct discrimination when $|\mathrm{DE}| > 0.05$; a mediator channel counted as negligible when $|\mathrm{IE}| < 0.005$; practical relevance when $|\mathrm{TV}| > 0.05$; a substantial spurious component when $|\mathrm{SE}| / |\mathrm{TV}| > 0.1$.


\vspace{1em}
\noindent\textbf{Answer consistency:} 1/5 (20.0\%). The expected answers are derived deterministically from the FairMind values alone, through threshold rules, and were added to the document after it was generated: the model never saw them.

\end{document}